\subsection{Definition of a graph}

DEFINITION 1. --- A graph \(^{(1)}\) is a quiver (S, F, o, t) equipped with an involution of F, denoted \(f \mapsto \overline{f}\) , without fixed point and such that \(t(\overline{f}) = o(f)\) for all \(f \in F\) .

The quiver (S, F, o, t) is said to be underlying this graph. For every arrow f of this quiver, we have, by definition, \(\overline{\overline{f}} = f\) , \(\overline{f} \neq f\) and \(t(\overline{f}) = o(f)\) . By applying this last relation to the arrow \(\overline{f}\) , one obtains the equality \(o(\overline{f}) = t(f)\) . The arrow \(\overline{f}\) is called the arrow opposite to f.

A pair of opposite arrows of the graph is called an edge of the graph. Each of the two arrows belonging to this pair is called an orientation of this edge. For this reason, the arrows of a graph are also called the oriented edges of the graph.

If G and G' are graphs, a graph morphism from G to G' is a quiver morphism \(\varphi\colon\mathrm{G}\to\mathrm{G}'\) such that \(\overline{\varphi(f)}=\varphi(\overline{f})\) for every arrow \(f\) of G.

Let G and G' be graphs; one says that G' is a subgraph of G if it is a subquiver and if the involution of Fl(G') is the restriction of that of Fl(G).

